%% ============================================================
%%  SECCIÓN 9 — Tratamiento de Datos
%% ============================================================

\section{Tratamiento de Datos y Resultados}
\label{sec:tratamiento}

En esta sección se presentan los cálculos numéricos efectuados para determinar las magnitudes físicas derivadas a partir de la ecuación de trabajo (Ec.~\eqref{eq:modelo_general}):

\begin{equation}
    y = f(x_1, x_2, \dots, x_n)
    \label{eq:trabajo_calculo}
\end{equation}

La propagación de la incertidumbre para variables independientes se realiza mediante el método de suma en cuadratura:
\begin{equation}
    \delta y = \sqrt{\sum_{i=1}^n \left( \frac{\partial f}{\partial x_i}\,\delta x_i \right)^2}
    \label{eq:propagacion_cuadratura}
\end{equation}

El error relativo porcentual respecto a los valores teóricos bibliográficos se evalúa mediante:
\begin{equation}
    E_r = \frac{|y - y_{\text{ref}}|}{y_{\text{ref}}} \times \qty{100}{\percent}
    \label{eq:error_relativo}
\end{equation}

\subsection{Cálculo para la Condición 1}

Sustitución numérica paso a paso:
\begin{align*}
    x_1 &= \qty{10,00 \pm 0,05}{\milli\meter} \\
    x_2 &= \qty{25,0 \pm 0,5}{\celsius}
\end{align*}

Resultado final:
\begin{equation}
    y_1 = (12{,}34 \pm 0{,}56) \times 10^{-3}\ \text{unidad}
    \label{eq:res_1}
\end{equation}
con un error relativo de $E_r = \qty{2,5}{\percent}$.
